% =====================================================================
% chapters/02-literature-review.tex
% PT: Capítulo 2 -- Revisão da literatura. Exemplo de: citações com
%     \cite e \textcite, várias fontes numa citação, citação com página,
%     fonte online e acrónimos no plural (\glspl).
% EN: Chapter 2 -- Literature review. Shows: citations with \cite and
%     \textcite, several sources in one citation, citation with a page
%     number, an online source and plural acronyms (\glspl).
% =====================================================================
\chapter{Literature Review}
\label{chap:review}

% PT: Parágrafo de abertura do capítulo.
% EN: Chapter opening paragraph.
This chapter shows how to cite sources. In a real dissertation, the
literature review presents the concepts the work builds on, summarises and
compares the related work, and ends by identifying the gap the dissertation
fills. Every source cited here must have an entry in
\texttt{bibliography/references.bib}; the reference list at the end of the
document is generated automatically and only contains the sources that are
cited.

% ---------------------------------------------------------------------
% PT: Exemplo em inglês
% EN: English example
% ---------------------------------------------------------------------
\section{English example}
\label{sec:review:en}

\subsection{Citing Sources}
\label{sec:review:citing}

% PT: \cite{chave} dá "[1]"; \textcite{chave} dá "Autor et al. [1]".
%     A chave é o identificador da entrada no ficheiro .bib.
% EN: \cite{key} gives "[1]"; \textcite{key} gives "Author et al. [1]".
%     The key is the identifier of the entry in the .bib file.
Deep learning methods learn representations of data with several levels of
abstraction \cite{LeCun2015}. \textcite{Goodfellow2016} give a comprehensive
introduction to the field, and \textcite{He2016} showed that residual
connections make it possible to train much deeper \glspl{cnn}.

% PT: Várias fontes de uma vez e citação de uma página concreta.
% EN: Several sources at once and citation of a specific page.
Several sources can be cited together \cite{LeCun2015,He2016}, and a
specific page can be added in square brackets \cite[p.~164]{Goodfellow2016}.
Online sources, such as software documentation, are cited in the same way
\cite{CTANbiblatexieee}.

\subsection{Summary of Related Work}
\label{sec:review:summary}

A literature review should do more than list papers. Group the works by
approach, compare them on the criteria that matter for your problem, and
close the chapter by stating clearly what is still missing. A comparison
table, like the one in \cref{tab:results}, is often a good way to summarise
the related work.

% ---------------------------------------------------------------------
% PT: Exemplo em português
% EN: Portuguese example
% ---------------------------------------------------------------------
\section{Exemplo em português}
\label{sec:review:pt}

\begin{otherlanguage}{portuguese}

Os métodos de aprendizagem profunda aprendem representações dos dados com
vários níveis de abstração \cite{LeCun2015}. Segundo
\textcite{Goodfellow2016}, estas representações são aprendidas diretamente a
partir dos dados, sem regras definidas à mão. \textcite{He2016} mostraram que
as ligações residuais permitem treinar \glspl{cnn} com muito mais camadas.

Para acrescentar uma nova fonte, crie uma entrada no ficheiro
\texttt{references.bib} com uma chave única (por exemplo,
\texttt{Silva2024}) e cite-a no texto com \texttt{\textbackslash
cite\{Silva2024\}}.

\end{otherlanguage}
